\BourbakiExerciseGroup

除非另有明确说明，本文所考虑的所有环均假定为交换环。

\begin{enumerate}
\def\labelenumi{\arabic{enumi})}
\tightlist
\item
  设 A 是一个全序环，B 是 A 的一个加法子群，且 B 对乘法封闭。
\end{enumerate}

\begin{enumerate}
\def\labelenumi{\alph{enumi})}
\item
  元素 \(x \in A\) 称为关于 B 是无穷大的，如果 \(|y| < |x|\) 对所有 \(y \in B\) 成立。证明 A 中关于 B 不是无穷大的元素所成的集 \(F(B)\) 是 A 的一个包含 B 的子环。
\item
  元素 \(x \in A\) 称为关于 B 是无穷小的，如果对 B 中所有 y \textgreater{} 0，都有 \(|x| < y\) 。如果对 B 中所有 y \textgreater{} 0，存在 \(z \in B\) 使得 0 \textless{} z \textless{} y，则 A 中关于 B 为无穷小的元素所成的集是 A 的一个子伪环 I(B)。此外，若对 B 中每一对满足 0 \textless{} y \textless{} z 的元素 y、z，存在 \(x \in B\) 使得 0 \textless{} xz \textless{} y，则 I(B) 是环 F(B) 的一个理想。
\end{enumerate}

\begin{enumerate}
\def\labelenumi{\arabic{enumi})}
\setcounter{enumi}{1}
\item
  设 A 为一个全序环，n 为 A 的幂零元集合（它是 A 的一个理想）。A 中不属于 n 的每个元素关于 n 都是无穷大的；商环 A/n 是全序的（VI，第 30 页，习题 4），并且是一个整环。
\item
  在域 \(\mathbf{Q}\) 中，设 \(\mathbf{P}\) 为由 0 与 \(\geqslant 1\) 的有理数（在通常的序下）所组成的集。证明 \(\mathbf{P}\) 满足公理 \((\mathrm{AP}_{\mathrm{I}})\) 、 \((\mathrm{AP}_{\mathrm{II}})\) 与 \((\mathrm{AP}_{\mathrm{III}})\) 。
\item
  \begin{enumerate}
  \def\labelenumii{\alph{enumii})}
  \tightlist
  \item
    设 K 是一个域，P 是 K 的一个满足条件 \((\mathrm{AP}_{\mathrm{I}})\) 、 \((\mathrm{AP}_{\mathrm{II}})\) 与 \((\mathrm{AP}_{\mathrm{III}})\) 的子集，并且 \(K^{2} \subset P\) （其中 \(K^{2}\) 表示 K 的元素的平方所成的集）；证明在由 P 定义的序中，x \textgreater{} 0 蕴含 \(x^{-1} > 0\) ；由此推出，集合 \(K_{+}^{*}\) （由 K 中 \textgreater{} 0 的元素组成）是 K 的乘法群的一个子群。
  \end{enumerate}
\end{enumerate}

\begin{enumerate}
\def\labelenumi{\alph{enumi})}
\setcounter{enumi}{1}
\item
  在域 \(\mathbf{Q}\) 中，集合 \(\mathbf{P}\) 若满足 \(a\) ) 的条件，则只能是 \(\geqslant 0\) 的有理数（在通常的序下）所成的集。
\item
  * 设 P 是 K 的一个子集，满足 \((\mathrm{AP}_{\mathrm{I}})\) 、 \((\mathrm{AP}_{\mathrm{II}})\) 与 \((\mathrm{AP}_{\mathrm{III}})\) ，使得 \(1 \in P\) ，并且使得在由 P 定义的序中，x \textgreater{} 0 蕴含 \(x^{-1} > 0\) 。证明对正的 y 与 z，\(y^{2} \geqslant z^{2}\) 蕴含 \(y \geqslant z\) 。由此推出，对任意的 y \textgreater{} 0，我们有
\end{enumerate}

\[
\frac {1}{2} (y + y ^ {- 1}) \geqslant 1 - n ^ {- 1}
\]

对每个整数 \(n > 0\) 成立（注意 \((y + y^{-1})^{2m} \geqslant \binom{2m}{m}\) 对每个整数 \(m > 0\) 成立）. 由此推出，若由 P 定义的序加法群结构是阿基米德的（VI，第 35 页，习题 31），则 \((y - z)^2 \geqslant 0\) 对 P 的每一对元素 \(y, z\) 成立；若 \(K'\) 是由 P 生成的 K 的子域，则 \(x^2 \geqslant 0\) 对所有 \(x \in K'\) 成立。 *

\(d) * \text{令 } K = R(X) \text{ 为关于一个不定元的有理函数域，其系数域为 } R; \text{ 令 } P \text{ 为由 0 以及有理函数 } u \in K \text{ 满足以下条件：} u(t) \text{ 有定义且 } >0 \text{ 对每个实数 } t. \text{ 证明 } P \text{ 满足条件 } c) \text{ 并生成 } K, \text{ 但存在元素 } v \in K \text{ 满足以下条件：} v^2 \notin P. *\)

\begin{enumerate}
\def\labelenumi{\arabic{enumi})}
\setcounter{enumi}{4}
\tightlist
\item
  设 A 是一个格序环。A 的一个凸理想（VI，第 30 页，习题 4）称为不可约的，如果它不是两个不同于自身的凸理想的交集。
\end{enumerate}

\begin{enumerate}
\def\labelenumi{\alph{enumi})}
\tightlist
\item
  证明 A 的不可约凸理想的交为 0（若 \(a \in \mathbf{A}\) 非零，考虑 A 的一个在所有不包含 \(a\) 的理想中极大的凸理想）。由此推出 A 同构于一个子环 \(\mathbf{A}'\) （某个乘积 \(\prod_{i} \mathbf{A}_i\) 的），使得每个 \(\mathbf{A}_i\) 都是格序的，
\end{enumerate}

\(\mathrm{pr}_{\iota}(\mathbf{A}^{\prime})=\mathbf{A}_{\iota}\) 对所有 \(\iota\) 成立，且理想 \(\{0\}\) 在 \(\mathbf{A}_{\iota}\) 中（作为凸理想）对所有 \(\iota\) 不可约。 b) 证明以下条件在 A 中等价：

α) \(|xy| = |x| \cdot |y|\) 对所有 x, y；

\(\beta) x \cdot \sup(y, z) = \sup(xy, xz)\) 对所有 \(x \geqslant 0\) , \(y, z\) ;

\(\gamma) x \cdot \inf(y, z) = \inf(xy, xz)\) 对所有 \(x \geqslant 0\) , \(y, z\) ;

\(\delta)\inf (y,z) = 0\) 蕴含 \(\inf (xy,z) = 0\) ，对所有 \(x\geqslant 0\) .

（为看出 \(\gamma\) ) 蕴含 \(\delta\) )，注意 \(xy \leqslant y \cdot \sup (x, 1)\) ).

当这些条件满足时，称A为强格序的。

\begin{enumerate}
\def\labelenumi{\alph{enumi})}
\setcounter{enumi}{2}
\tightlist
\item
  证明一个环 \(\mathbf{A}\) 是强格序的当且仅当它同构于一个乘积 \(\prod_{i} \mathbf{A}_{i}\) 的子环，其中每个因子是全序环。（利用 \(a\) ），归结为证明：若
\end{enumerate}

\{0\} 在强格序环 A 中不可约，则 A 是全序的；为此，注意若 B 是强格序环，则对所有 \(b \in B\) ，集合 \(m\) 由这样的元素 \(x \in B\) 组成：\(\inf(|x|, |b|) = 0\) ；集合 \(n\) 由这样的元素 \(y \in B\) 组成：\(|y| \leqslant |zb|\) 对某个 \(z \in B\) 成立；它们是 B 的两个凸理想，且满足 \(m \cap n = \{0\}\) .

\(d\) ) 证明在强格序环 \(\mathbf{A}\) 中，关系 \(\inf (x,y) = 0\) 蕴含 \(xy = 0\) ；对所有 \(z\in \mathbf{A}\) 有 \(z^2\geqslant 0\) ；对所有 \(x,y\in \mathbf{A}\) 有 \(xy\leqslant \sup (x^2,y^2)\)

\begin{enumerate}
\def\labelenumi{\alph{enumi})}
\setcounter{enumi}{4}
\tightlist
\item
  设 A 是一个没有 \textgreater0 幂零元素的格序环；证明：若 \(\inf(x, y) = 0\) 蕴含 xy = 0，则 A 是强格序的（验证条件 \(\delta\) ) ，即 b) 中所述）。
\end{enumerate}

\begin{enumerate}
\def\labelenumi{\arabic{enumi})}
\setcounter{enumi}{5}
\tightlist
\item
  \begin{enumerate}
  \def\labelenumii{\alph{enumii})}
  \tightlist
  \item
    设 K 是一个格序域。证明以下条件等价：
  \end{enumerate}
\end{enumerate}

\(\alpha) x^{2} \geqslant 0\) 对所有 \(x \in \mathbf{K}\) ;

\(\beta) x > 0\) 蕴含 \(x^{-1} > 0\) ;

γ) K 是强格序的（习题 5）；

δ) K 是全序的。

（为看出 \(\beta\) ) 蕴含 \(\underline{\delta}\) ），注意 \(\beta\) 蕴含 \(xy \leqslant x^2 + y^2\) ，对 \(x, y \in K\) ).

\begin{enumerate}
\def\labelenumi{\alph{enumi})}
\setcounter{enumi}{1}
\item
  设 K 为域 \(\mathbf{Q}(\sqrt{2})\) ，它是由 \(\sqrt{2}\) 添加到有理数域而得到的；作为加法群，K 与 \(Q \times Q\) 通过双射 \(x + y \sqrt{2} \mapsto (x, y)\) 等同；证明：若将 \(Q \times Q\) 上的乘积序（其中 Q 具有通常的序）经由该双射转移到 K 上，则 K 是格序的，但不是强格序的。
\item
  域 \(\mathbf{K} = \mathbf{Q}(\mathbf{X})\) （Q 上关于单个未定元 X 的有理函数域）由它的平方所成的集 \(K^{2}\) 生成。证明：K 中元素的平方和所成的集 P 在 K 上定义了一个与环结构相容的非格序，并且使得 u \textgreater{} 0 蕴含 \(u^{-1} > 0\) .
\end{enumerate}

\begin{enumerate}
\def\labelenumi{\arabic{enumi})}
\setcounter{enumi}{6}
\item
  特征为 \(\neq 2\) 的交换域 K 中，每个元素都是平方和，当且仅当 -1 是平方和（注意 K 的每个元素都是两个平方之差）。
\item
  设 A 是特征为 \(\neq 2\) 的交换域 K 的一个非空子集；若 A 中没有元素是 K 中平方和，则称 K 是 A-可序的。我们称 K 是可序的，如果存在 K 的一个全序，且该全序与其环结构相容。
\end{enumerate}

\begin{enumerate}
\def\labelenumi{\alph{enumi})}
\item
  证明若 K 是 A-可序的，则它是可序的（习题 7），从而特征为 0。
\item
  证明 A-可序域 K 的纯扩张和奇数次代数扩张是 A-可序的（仿照 VI 第 23 页命题 3 论证）。
\item
  设 K 是 A-可序域，b 是 K 中一个元素，且 b 不具有形式 ca - d，其中 \(a \in A\) 而 c 和 d 是 K 中平方和。证明域 \(\mathrm{K}(\sqrt{b})\) 是 A-可序的。
\item
  A-可序域 K 称为极大的，如果 K 没有真代数扩张是 A-可序的。证明极大 A-可序域 K 具有以下性质：1. K 是毕达哥拉斯域，换言之每个平方和都是平方；2. A 中没有元素是平方；3. K 中每个非平方元素都具有形式 \(c^2 a - d^2\) ，其中 \(a \in \mathbf{A}\) ；4. K{[}X{]} 中每个奇数次多项式在 K 中至少有一个根（利用 b）与 c））。e) 证明若 K 满足 d) 的条件，则 K 上每个代数元在 K 上的次数都形如 \(2^q\) （对某个 q）（仿照 VI 第 26 页命题 8，对所考虑元素的次数中 2 的指数作归纳）。另一方面，证明 K 的任何二次扩张都不是 A-可序的。由此推出 K 是极大 A-可序域（利用伽罗瓦理论和 I 第 77 页命题 12）。
\item
  设 K 是 A-可序域，且 \(\Omega\) 是 K 的代数闭扩张。证明存在一个包含于 \(\Omega\) 且包含 K 的极大 A-可序域。
\end{enumerate}

\begin{enumerate}
\def\labelenumi{\arabic{enumi})}
\setcounter{enumi}{8}
\tightlist
\item
  设 \(q > 0\) 为一个非平方自然数；令 \(A\) 表示集合 \(\{-1, \sqrt{q}\}\) （在域 \(K = Q(\sqrt{q})\) 中）；证明 \(K\) 是 \(A\) -可序的。证明存在一个扩张 \(E\) （ \(K\) 的），使得 \(E\) 是毕达哥拉斯的、可序的，并且 \(E\) 上每个奇数次多项式在 \(E\) 中有一个根，但 \(E\) 不具有极大序域的结构（考虑 \(A\) -可序的 \(K\) 的一个极大扩张）.
\end{enumerate}

¶ 10) a) 设 K 是一个可序域，E 是 K 的一个伽罗瓦扩张。证明：要么 E 是可序的，要么存在 K 的一个包含于 E 的可序代数扩张 F，使得 E 是 F 的二次扩张（使用 V 第 70 页定理 6）。

\begin{enumerate}
\def\labelenumi{\alph{enumi})}
\setcounter{enumi}{1}
\tightlist
\item
  证明多项式 \(X^{4} + 2\) 在 Q 上不可约，并且在 Q 上添加该多项式的一个根所得到的 Q 的 4 次扩张 K 不包含除 Q 以外的任何可序子域（利用伽罗瓦理论找出 K 的所有子域）。
\end{enumerate}

\begin{enumerate}
\def\labelenumi{\arabic{enumi})}
\setcounter{enumi}{10}
\tightlist
\item
  设 K 是一个序域，G 是 K 的一个子域。a) 证明：\(x \neq 0\) 关于 G 是无穷大的，当且仅当 \(x^{-1}\) 关于 G 是无穷小的。我们说 K 与 G 可比较，如果 K 中没有关于 G 无穷大的元素（从而也没有 K 的非零元素关于 G 是无穷小的）。K 与它的素子域 Q 可比较的充要条件是 K 是阿基米德的（VI，第～35 页，习题 31），* 从而 K 同构于 R 的一个子域（VI，第～36 页，习题 33c））。
\end{enumerate}

\begin{enumerate}
\def\labelenumi{\alph{enumi})}
\setcounter{enumi}{1}
\item
  证明在由 K 中关于 G 不是无穷大的元素组成的环 \(F(G)\) 中，关于 G 为无穷小的元素组成的集 \(I(G)\) 是一个极大理想；此外，商域 \(K(G) = F(G)/I(G)\) 上由 \(F(G)\) 的序诱导的序与环结构相容，并且是全序的。
\item
  证明模 \(I(G)\) 的一个等价类至多包含 G 的一个元素；由此推出从 G 到 \(K(G)\) 的自然映射是序域 G 到子域 \(G'\) （ \(K(G)\) 的）上的一个同构，并且 \(K(G)\) 与 \(G'\) 可比较。
\end{enumerate}

\begin{enumerate}
\def\labelenumi{\arabic{enumi})}
\setcounter{enumi}{11}
\item
  设 K 为一个极大序域，f 为 K 上的一个多项式，并设 a 和 b 为 f 的两个根，满足 a \textless{} b，且使 f 在 a 与 b 之间没有根。证明：若 g 是 K 上的有理函数，其分母在 \(a \leqslant x \leqslant b\) 上不消失，则方程 \(f(x)g(x) + f'(x) = 0\) 在 \([a, b]\) 中有奇数个解（其中每个解都按重数计算；使用 VI 第 25 页命题 5）。由此推出，如果 h 是 K 上的有理函数，以 a 和 b 为根，且其分母在 \([a, b]\) 中不消失，则方程 \(h'(x) = 0\) 在 \([a, b]\) 中至少有一个根。
\item
  设 K 是一个极大序域，h 是 K 上的一个有理函数，且 \([a, b]\) 是 h 在其上有定义的闭区间。证明存在 \(c \in ]a, b[\) 使得 \(h(b) - h(a) = (b - a)h'(c)\) （利用习题 12）。由此推出，h 在 \([a, b]\) 上是递增函数，当且仅当在此区间内 \(h'(x) \geqslant 0\) （要看出该条件是必要的，可用 \(h'(x) = 0\) 的根分割该区间）.
\item
  \begin{enumerate}
  \def\labelenumii{\alph{enumii})}
  \tightlist
  \item
    设 K 为一个极大序域，E 为 K 的子域，f 为 E 上的多项式；证明 f 在 K 中的所有根都属于 F(E)（习题 11，b））（利用 VI 第 24 页命题 4）。由此推出，若 G 为 K 的子域且 E ⊂ K 是与 G 可比较的扩张（习题 11），则 K 中在 E 上代数的元素之集是一个与 G 可比较的极大序域。
  \end{enumerate}
\end{enumerate}

\begin{enumerate}
\def\labelenumi{\alph{enumi})}
\setcounter{enumi}{1}
\item
  由a)推出，在a)的假设下，域K(G)（习题11，b)）是一个极大序域。
\item
  设 \(f\) 是次数 \(\geqslant 1\) 的 \(G\) 上的多项式。证明 \(f(t)\) 关于 \(G\) 是无穷大的，当且仅当 \(t\) 关于 \(G\) 是无穷大的；元素 \(f(t)\) 关于 \(G\) 是无穷小的，当且仅当 \(t\) 模 \(I(G)\) 同余于 \(f\) 在 \(K\) 中的一个根（把 \(K\) 用 \(f\) 与 \(f'\) 在 \(K\) 中的根分割成区间，并应用习题 13；注意若 \(x < t\) 且 \(x \in K\) 不模 \(I(G)\) 同余于 \(t\) ，则存在 \(y \in K\) 使得 \(x < y < t\) 且 \(y - x \in G\) ).
\end{enumerate}

\begin{enumerate}
\def\labelenumi{\arabic{enumi})}
\setcounter{enumi}{14}
\tightlist
\item
  设 E 是一个序域，K 是 E 的子域，使得 E 在 K 上是代数的，并且 R 是 K 的一个极大序扩张。
\end{enumerate}

\begin{enumerate}
\def\labelenumi{\alph{enumi})}
\item
  在 E - K 的元素中，设 \(x_0\) 是其在 K 上的极小多项式 \(f\) 的次数尽可能小的一个。证明 \(f'(T)\) 是 R{[}T{]} 中形如 \((T - a_i)^2 + c_i^2\) ( \(a_i, c_i \in \mathbb{R}\) ) 的因子与一次因子 \(T - b_j\) （其中 \(b_j\) 是 K 中两两不同的元素）的乘积。由此推出，存在 R 中的一个元素 \(y_0\) ，使得 \(f(y_0) = 0\) ，且 \(x_0 - z\) 与 \(y_0 - z\) 分别在 E 与 R 中对所有 \(z \in \mathbb{K}\) 具有相同的符号（参见 VI，第 43 页，习题 26，c））。
\item
  由 a) 推出，存在一个从序域 E 到 R 的某个子域上的 K-同构。（首先证明对 E 的子域 \(\mathrm{K}(x_{0})\) 存在这样的同构；对任意多项式 \(g \in \mathrm{K}[\mathrm{T}]\) （满足 \(\deg(g) < \deg(f)\) ），用 a) 中对 \(f'\) 用过的同样论证来求 \(g(x_{0})\) 与 \(g(y_{0})\) 的符号；然后应用佐恩引理）。
\item
  证明：若R和R'是K的两个极大序代数扩张，则存在恰好一个从R到R'的K-同构（应用b)）。
\end{enumerate}

\begin{enumerate}
\def\labelenumi{\arabic{enumi})}
\setcounter{enumi}{15}
\tightlist
\item
  设K是一个极大序域，并设
\end{enumerate}

\[
g (\mathrm{T}) = a _ {0} + a _ {1} \mathrm{T} + \dots + a _ {n} \mathrm{T} ^ {n}
\]

是 K{[}T{]} 中的一个非零多项式，其根全部位于 K 中。证明对于任意多项式 \(f \in K[T]\) ，多项式

\[
a _ {0} f (\mathrm{T}) + a _ {1} f ^ {\prime} (\mathrm{T}) + a _ {2} f ^ {\prime \prime} (\mathrm{T}) + \dots + a _ {n} f ^ {(n)} (\mathrm{T})
\]

的不属于 K 的根（按重数计）的个数至多等于 f 的不属于 K 的根（按重数计）的个数（对 n = 1 使用习题 12，然后归纳进行）。由此推出多项式

\[
a _ {0} + \frac {a _ {1}}{1 !} \mathrm{T} + \frac {a _ {2}}{2 !} \mathrm{T} ^ {2} + \dots + \frac {a _ {n}}{n !} \mathrm{T} ^ {n}
\]

的所有根都在K中。

\begin{enumerate}
\def\labelenumi{\arabic{enumi})}
\setcounter{enumi}{16}
\tightlist
\item
  设 K 是一个极大序域，g 是 K{[}T{]} 中的非零多项式，其根均在 K 中，且不属于 K 的区间 \([0, n]\) 。证明对于每一个多项式 \(f(T) = a_{0} + a_{1}T + \cdots + a_{n}T^{n}\) （属于 K{[}T{]}），多项式
\end{enumerate}

\[
a _ {0} g (0) + a _ {1} g (1) \mathrm{T} + a _ {2} g (2) \mathrm{T} ^ {2} + \dots a _ {n} g (n) \mathrm{T} ^ {n}
\]

的不属于 K 的根（按重数计）的个数至多等于 f 的不属于 K 的根（按重数计）的个数（方法与习题 16 相同）。

\begin{enumerate}
\def\labelenumi{\arabic{enumi})}
\setcounter{enumi}{17}
\item
  设 K 为一个极大序域，f 为 K{[}T{]} 中次数为 n 的多项式，其所有根都在 K 中。证明对 K 中所有 \(c \neq 0\) ，多项式 \(f^{2} + cf'\) 在 K 中按重数计至少有 n - 1 个且至多有 \(n + 1\) 个根（使用习题 13）。
\item
  设 K 是一个极大序域，\(f\) 是 K{[}T{]} 中的一个非零多项式。使得对每个多项式 \(g \in \mathbf{K}[\mathbf{T}]\) ，\(fg + g'\) 的不属于 K 的根（按重数计）的个数至多等于 \(g\) 的不属于 K 的根（按重数计）的个数，其必要且充分条件是 \(f(\mathbf{T}) = a - b\mathbf{T}\) ，其中 \(b \geqslant 0\) （在 K 中）。（要看出该条件是充分的，使用习题 12；要证明它是必要的，取 g = 1 和 g = f）。
\item
  设 \(\mathbf{K}\) 是一个极大序域，且
\end{enumerate}

\[
f (\mathrm{T}) = a _ {0} + \binom {n} {1} a _ {1} \mathrm{T} + \binom {n} {2} a _ {2} \mathrm{T} ^ {2} + \dots + a _ {n} \mathrm{T} ^ {n}
\]

是 K{[}T{]} 中的一个多项式。对于所有满足 \(0 \leqslant p < p + q \leqslant n\) 的整数 p, q，多项式

\[
a _ {p} + \binom{q}{1} a _ {p + 1} \mathrm{T} + \binom{q}{2} a _ {p + 2} \mathrm{T} ^ {2} + \dots + a _ {p + q} \mathrm{T} ^ {q}
\]

（按重数计）不属于 K 的根的个数至多等于 f 的（按重数计）不属于 K 的根的个数（利用习题 12）。

由此推出，如果 \(b_{1}, \ldots, b_{n}\) 是 K 中 n 个两两不同的 \textgreater0 的元素，且若

\[
\left(\mathrm{T} + b _ {1}\right) \dots \left(\mathrm{T} + b _ {n}\right) = \mathrm{T} ^ {n} + \binom {n} {1} m _ {1} \mathrm{T} ^ {n - 1} + \dots + m _ {n}
\]

则

\[
m _ {k} ^ {1 / k} > m _ {k + 1} ^ {1 / (k + 1)} \quad \text { for } \quad 1 \leqslant k \leqslant n - 1.
\]

\begin{enumerate}
\def\labelenumi{\arabic{enumi})}
\setcounter{enumi}{20}
\tightlist
\item
  设 \((a_{i})_{1 \leqslant i \leqslant n}\) 为序域 K 中 n 个元素组成的有限序列，且不全为零；设 \((a_{i_{k}})_{1 \leqslant k \leqslant p}\) 为 \((a_{i})\) 的由非零 \(a_{i}\) 组成的子序列 ( \(i_{1} < i_{2} < \cdots < i_{p}\) ）；则指标 \(k \leqslant p - 1\) （使得 \(a_{i_{k}}\) 与 \(a_{i_{k+1}}\) 具有相反符号者）的个数称为序列 \((a_{i})\) 的变号数。
\end{enumerate}

设 K 为一个极大序域，f 为 K{[}T{]} 中次数为 n \textgreater{} 0 的非零多项式；对所有 \(a \in K\) ，令 \(w(a)\) 表示序列 \((f^{(i)}(a))_{0 \leq i \leq n}\) 的变号数。若 v 是 f 在区间 \([a, b]\) (a \textless{} b) 内按重数计的根的个数，证明 \(v \leq w(a) - w(b)\) 且 \(w(a) - w(b) - v\) 是偶数（“布当–傅里叶法则”；用 f 的根分割区间 \([a, b]\) ，并计算当 x 经过其中一个根时 \(w(x)\) 的变号数的变化）。

由此推出，如果 \(f(T) = a_{0} + a_{1}T + \cdots + a_{n}T^{n}\) ，则 f 的按重数计的 \textgreater{} 0 的根的个数至多等于序列 \((a_{i})_{0 \leq i \leq n}\) 的变号数，并且这两个数之差是偶数（“笛卡儿法则”）。

\begin{enumerate}
\def\labelenumi{\arabic{enumi})}
\setcounter{enumi}{21}
\tightlist
\item
  设K是一个极大序域，并设
\end{enumerate}

\[
f (\mathrm{T}) = a _ {0} + a _ {1} \mathrm{T} + \dots + a _ {n} \mathrm{T} ^ {n}
\]

是 K{[}T{]} 中的多项式，满足 \(a_0 \neq 0 \neq a_n\) ，且

\[
a _ {p} = a _ {p + 1} = \dots = a _ {p + 2 m - 1} = 0 \quad (1 \leqslant p <   p + 2 m - 1 <   n).
\]

证明 \(f\) 按重数计至多有 \(n - 2m\) 个根在 \(\mathbf{K}\) 中（应用笛卡儿法则）。

由此推出，如果 \(g(T) = 1 + c_{1}T + \cdots + c_{n}T^{n}\) 的所有根都在 K 中，并且 \(h(T) = 1 + b_{1}T + \cdots + b_{2m}T^{2m}\) 是由前 \(2m + 1\) 项组成的多项式，这些项取自形式幂级数 \(1/g \in K[[T]]\) ，则 h 在 K 中没有根。

\begin{enumerate}
\def\labelenumi{\arabic{enumi})}
\setcounter{enumi}{22}
\tightlist
\item
  设 K 是一个极大序域，并设 \(a_{1}, \ldots, a_{n}, b_{1}, \ldots, b_{n}\) 是 K 的 2n 个元素，使得 \(\sum a_{i} \neq 0\) 且 \(b_{1} < b_{2} < \cdots < b_{n}\) 。证明对于任意整数 \(m \geqslant 1\) ，多项式在 K 中的根（按重数计）的个数 v
\end{enumerate}

\[
f (\mathrm{T}) = a _ {1} \left(\mathrm{T} - b _ {1}\right) ^ {m} + a _ {2} \left(\mathrm{T} - b _ {2}\right) ^ {m} + \dots + a _ {n} \left(\mathrm{T} - b _ {n}\right) ^ {m}
\]

至多等于序列的变号数w，

\[
\left(a _ {1}, a _ {2}, \dots , a _ {n}, (- 1) ^ {m} a _ {1}\right),
\]

且差w - v为偶数。（通过对n归纳证明，将习题13（VI，第40页）应用于形如 \(f(\mathrm{T})/(\mathrm{T}-c)^{m}\) .)

\begin{enumerate}
\def\labelenumi{\arabic{enumi})}
\setcounter{enumi}{23}
\item
  设 K 为一个极大序域；考虑有理函数域 K(X) 上的一个序，使其成为 K 的序扩张。证明该序由集合 A 决定，其中 A 为 \(x \in K\) 使得 x \textless{} X 的集合（利用 VI 第 28 页命题 9 证明 K 上每个多项式 f 的符号由 A 决定）。反之，证明对每个满足 \(y \in A\) 当 \(y \leqslant x\) 中，且 \(x \in A\) 时的集合 A ⊂ K，都存在 K(X) 的一个序，使其成为 K 的序扩张，并且 A 恰为满足 \(x \in K\) 使得 x \textless{} X 的集合（用同样方法）。当 A 有最大元，或 CA 有最小元，或 A = K 或 A = ∅ 时，对应的 K(X) 的序使得 K(X) 与 K 不可比较。另一方面，若 A 与 CA 均非空，且 A 无最大元、CA 无最小元，则 K(X) 与 K 可比较。
\item
  利用习题 24 找出一个域 E 的例子，使其具有多个序域结构，且其中任何一个都不能通过 E 的自同构由另一个诱导得到。由此找出一个域的例子，它带有与其环结构相容的序，但不是格序（考虑 E × E 中的对角线，并利用 VI 第 38 页习题 5）。
\item
  \begin{enumerate}
  \def\labelenumii{\alph{enumii})}
  \tightlist
  \item
    若 K 是阿基米德序域（VI 第 40 页习题 11 a)），且 x 与 y 是 K 中满足 x \textless{} y，证明存在 \(r \in Q\) 使得 x \textless{} r \textless{} y. * 推导出存在唯一一个与 K 同构的 R 的序子域。 *
  \end{enumerate}
\end{enumerate}

\begin{enumerate}
\def\labelenumi{\alph{enumi})}
\setcounter{enumi}{1}
\tightlist
\item
  考虑域 \(\mathbf{K} = \mathbf{Q}(\mathbf{X})\) 的序，其中 \(\mathbf{X}\) 关于 \(\mathbf{Q}\) 是无穷大的（习题 24）。证明 \(\mathbf{K}\) 与子域 \(\mathbf{Q}(\mathbf{X}^2)\) 可比较，并给出两个元素 \(x, y\) （属于 \(\mathbf{K}\) ）的例子，使得 \(x < y\) 且 \(\mathbf{Q}(\mathbf{X}^2)\) 中没有元素位于 \(x\) 与 \(y\) 之间。c) 设 \(\mathbf{K}\) 为在 \(b\) 中定义的序域；证明多项式
\end{enumerate}

\[
f (\mathbf {Y}) = (\mathbf {Y} ^ {2} - \mathbf {X}) (\mathbf {Y} ^ {2} - 4 \mathbf {X}) - 1
\]

在 K 上不可约，它在 K 的每个极大序扩张中都有根，并且 \(f(a) > 0\) 对所有 \(a \in K\) .

\begin{enumerate}
\def\labelenumi{\arabic{enumi})}
\setcounter{enumi}{26}
\tightlist
\item
  设 K 是一个序域，设 E 是 K 的一个纯超越扩张，并设 \((\mathbf{X}_{i})_{i\in I}\) 是 E 的一个超越基（V，第 109 页）。
\end{enumerate}

\begin{enumerate}
\def\labelenumi{\alph{enumi})}
\item
  * 若 K 是阿基米德的，则 E 具有 K 的序扩张结构，且 E 与 K 可比较，当且仅当 I 的基数至多等于 R 在 K 上的超越基的基数（其中 K 被视为 R 的序子域，参见习题 26 a)）；此时，此类结构的集合与从 I 到 R 的、使 \(f(I)\) 在 K 上代数无关的单射映射所成的集合成一一对应。*
\item
  如果K不是阿基米德的，那么在E上总是存在（至少）一种结构，使得E成为K的序扩张，并且E与K是可比较的（当I只包含一个元素时，使用习题24并注意到Q在K中没有上确界；通过佐恩引理过渡到一般情形）。
\end{enumerate}

* 28) a) 设 K 是 R 的一个子域，并且设 \(\theta\) 是在 K 上代数的实数。证明 K(θ) 上作为 K 的序扩张的结构的个数等于 \(\theta\) 的实共轭元的个数（利用习题 14，a）（VI，第 40 页）和 26，a））。

\begin{enumerate}
\def\labelenumi{\alph{enumi})}
\setcounter{enumi}{1}
\tightlist
\item
  设 K 为 R 的一个子域，带有 n 个不同的作为 Q 的序扩张的结构，且均为阿基米德的；设 \(P_{i}\) ( \(1 \leqslant i \leqslant n\) ) 为这些序下的正元素集合。证明存在一个由 K 中 n 个元素组成的族 \(b_{i}\) ( \(1 \leqslant i \leqslant n\) )，使得 \(b_{i} \in P_{i}\) 对所有 i 成立，且 \(b_{i} \notin P_{k}\) 对所有 \(k \neq i\) 成立。（对于每一对不同的指标 i, k，设 \(a_{ik} \in P_{k}\) 使得 \(a_{ik} \notin P_{i}\) ；考虑元素
\end{enumerate}

\[
c _ {i} = \sum_ {k \neq i} \left(\frac {1 + a _ {i k}}{1 - a _ {i k}}\right) ^ {2 r}
\]

，其中 \(r\) 是一个足够大的整数。）

\begin{enumerate}
\def\labelenumi{\alph{enumi})}
\setcounter{enumi}{2}
\tightlist
\item
  设 \(\mathbf{K} = \mathbf{Q}(\theta)\) 为 \(\mathbf{Q}\) 的一个包含于 \(\mathbf{R}\) 的代数扩张，并设 \(n\) 为 \(\theta\) 的实共轭元的个数。设 \(\mathbf{C}\) 为 \(\mathbf{K}\) 中元素的平方和所构成的集合；证明 \(\mathbf{C}^* = \mathbf{C} \cap \mathbf{K}^*\) 是 \(\mathbf{K}^*\) 的一个子群，且 \((\mathbf{K}^*: \mathbf{C}^*) = 2^n\) 。（使用 \(b\) ）以及 VI 第 23 页推论 2。）*
\end{enumerate}

\begin{enumerate}
\def\labelenumi{\arabic{enumi})}
\setcounter{enumi}{28}
\item
  设 K 为一个极大序域，G 为 K 的子域。证明：包含于 K 中且与 G 可比较的 G 的扩张之集是归纳有序的；若 \(E_{0}\) 是该集合的一个极大元素，证明 \(E_{0}\) 同构于 VI 第 40 页习题 11, b) 中定义的域 \(\mathrm{K}(\mathrm{G})\) （证明从 \(\mathrm{F}(\mathrm{G})\) 到 \(\mathrm{K}(\mathrm{G})\) 上的自然映射把 \(E_{0}\) 映到 \(\mathrm{K}(\mathrm{G})\) 上：首先利用 VI 第 40 页习题 14, a) 证明 \(E_{0}\) 是一个极大序域，进而利用习题 24 证明 \(\mathrm{K}(\mathrm{G})\) 中没有任何元素在 \(E_{0}\) 的像上是超越的）。
\item
  \begin{enumerate}
  \def\labelenumii{\alph{enumii})}
  \tightlist
  \item
    设 K 是一个极大序域，并设 m 和 M 是 K 中满足 m \textless{} M 的两个元素。证明每个多项式 \(f \in K[X]\) （在区间 \((m, M)\) 内为正的）都是形如 \((\alpha X + \beta) g^{2}\) 的多项式之和，其中 \(g \in K[X]\) ，且 \(\alpha X + \beta\) 在 \((m, M)\) 内为正。（化归为一次或二次多项式的情形，并利用公式
  \end{enumerate}
\end{enumerate}

\[
\begin{array}{l} (\mathbf {X} - a) (\mathbf {X} - b) = (\mathbf {X} - b) ^ {2} + (b - a) (\mathbf {X} - b) \\ (\mathbf {X} - a) (b - \mathbf {X}) = ((\mathbf {X} - a) (b - \mathbf {X}) ^ {2} + (b - \mathbf {X}) (\mathbf {X} - a) ^ {2}) / (b - a) \end{array}
\]

对 \(a < b\) .)

\begin{enumerate}
\def\labelenumi{\alph{enumi})}
\setcounter{enumi}{1}
\tightlist
\item
  证明当 K 为任意序域时，a) 的结果不再必然成立（注意一个多项式在 K 中可以为 \textgreater0，但在 K 的一个极大序扩张中不为正；参见习题 26，c)）。
\end{enumerate}

\begin{enumerate}
\def\labelenumi{\arabic{enumi})}
\setcounter{enumi}{30}
\tightlist
\item
  \begin{enumerate}
  \def\labelenumii{\alph{enumii})}
  \tightlist
  \item
    设 \(\mathbf{K}\) 是一个域，其代数闭包 \(\mathbf{E}\) 是 \(\mathbf{K}\) 的素数次 \(q\) 扩张。证明 \(\mathbf{K}\) 是完全域（V，第 7 页）。
  \end{enumerate}
\end{enumerate}

\begin{enumerate}
\def\labelenumi{\alph{enumi})}
\setcounter{enumi}{1}
\item
  证明 q 不同于 K 的特征（V，第 162 页，习题 9）。
\item
  证明K包含q次单位根，且E是形如 \(X^{q}-a\) 的不可约多项式的分裂域；由此推出q=2（否则，利用习题5，V，第161页推出 \(X^{q^{2}}-a\) 将是不可约的）。同时证明 -a 在 K 中是平方（见 V，前引），-1 在 K 中不是平方，并且 \(E = K(i)\) ( \(i^{2} = -1\) ).
\item
  现在假设 K 的代数闭包 E 在 K 上具有任意有限次数 \(\neq1\) 。证明 \(i\notin K\) 且 \(\mathrm{E}=\mathrm{K}(i)\) （若 \(\mathrm{E}\neq\mathrm{K}(i)\) ，用伽罗瓦理论证明会存在一个域 F 使得 \(\mathrm{K}(i)\subset\mathrm{F}\subset\mathrm{E}\) ，并且 E 在 F 上的次数为素数；然后应用 c)）。由此推出 K（在某个适当的序下）是一个极大序域（用归纳法证明 K 中任意 n 个平方的和是 K 中的平方；为了证明 \(a^{2}+b^{2}\) 是平方，考虑平方根 \(x+iy\) （ \(a+ib\) 在 \(\mathrm{K}(i)\) 中的）).
\end{enumerate}

\begin{enumerate}
\def\labelenumi{\arabic{enumi})}
\setcounter{enumi}{31}
\tightlist
\item
  设 A 是一个特征为 0 的代数闭域。
\end{enumerate}

\begin{enumerate}
\def\labelenumi{\alph{enumi})}
\item
  在群 \(\operatorname{Aut}(\mathbf{A})\) （ \(\mathbf{A}\) 的自同构群）中，仅有的有限阶元素是单位元和二阶元素（ \(\mathbf{A}\) 的对合）；后者与极大序子域 \(\mathbf{E}\) （ \(\mathbf{A}\) 中满足 \([\mathbf{A}:\mathbf{E}] = 2\) 者）一一对应，即 \(\mathbf{A} = \mathbf{E}(i)\) （参见习题 31）。若 \(\sigma\) 是 \(\mathbf{A}\) 的一个对合，且 \(\mathbf{E}\) 是 \(\sigma\) 的固定子域，则群 \(\operatorname{Aut}(\mathbf{E})\) （ \(\mathbf{E}\) 的自同构群）同构于商 \(Z(\sigma) / \{e, \sigma\}\) ，其中 \(Z(\sigma)\) 是 \(\sigma\) 在 \(\operatorname{Aut}(\mathbf{A})\) 中的中心化子。
\item
  如果 A 在其素域 Q 上是代数的，那么 A 的任意两个对合在 Aut(A) 中共轭，且对合 \(\sigma\) 的中心化子是 \(\{e, \sigma\}\) ；群 Aut(A) 与 A 的对合集合具有连续统的基数，且 Aut(A) 的中心是 \(\{e\}\) .
\item
  若 A 在 Q 上具有有限的超越次数，证明中心化子 \(Z(\sigma)\) （其中 \(\sigma\) 是 A 的任意一个对合）是可数的（注意 \(\sigma\) 的固定子域 E 是可数的，并且 E 的任何自同构，只要固定 E 在 Q 上的某个超越基的元素，就是恒等映射）。如果 deg. \(tr_{Q}A = n\) ，证明至少存在 \(n + 1\) 个 A 的对合的共轭类（若 \(E \subset A\) 是极大可序的且 \([A : E] = 2\) ，考虑 E 中由 Aut(E) 固定的元素所构成的子域在 Q 上的超越次数，并利用习题 24）。
\item
  若 A 在 Q 上具有无限超越次数 m，证明存在 A 的对合 \(\sigma\) ，使得 \(\operatorname{Card}(Z(\sigma)) = \operatorname{Card}(\operatorname{Aut}(A)) = 2^{m}\) （利用习题 24 证明存在极大可序域 E，使得 deg. \(tr_{Q}E = m\) ，使得对于 E 在 Q 上的任意超越基 B 的任意子集 M，存在 E 的一个自同构 \(u_{M}\) 满足： \(u_{M}(x) \neq x\) （对 \(x \in M\) ），且 \(u_{M}(x) = x\) （对 \(x \in B - M\) ）).
\item
  如果 \(\deg \cdot tr_{Q}A = 1\) ，证明 A 的包含无限中心化子的对合只有一个共轭类。（若 E 是 Q 的超越次数为 1 的极大可序扩张，且 \(\operatorname{Aut}(E)\) 非平凡，则利用习题 24 证明 E 不能与 Q 可比；若 \(t \in E\) 关于 Q 是无限大的，则证明对 E 的任何自同构 u， \(u(t)\) 也是无限大的，而反之，若 \(t'\) 关于 Q 是无限大的，则存在 E 的自同构 u 使得 \(u(t) = t'\) ，利用 VI 第 40 页习题 15。）
\item
  设 u 是 A 的一个自同构，且 u 与 A 的所有对合交换。证明若 \(a \in A\) 在 Q 上是超越的，则 a 和 \(u(a)\) 在 Q 上代数相关。（考虑 A 中两个元素 b, c，使得 \(b^{2} = a\) , \(c^{2} = -u(a)\) 并观察到存在一个极大的可序扩张 \(E \subset A\) 使得 \([A : E] = 2\) 且 b, c 属于 E；注意到 \(u(E) = E\) 便得出矛盾。）
\end{enumerate}

\(g\) ）由 \(f\) ) 推出 \(u(a) = a\) 对每个元素 \(a \in \mathbf{A}\) （它在 \(\mathbf{Q}\) 上是超越的）成立（考虑 \(\mathbf{Q}(a)\) 的一个极大序扩张 E，它包含于 A 中，在 \(\mathbf{Q}(a)\) 上是代数的，且与 \(\mathbf{Q}\) 可比，利用习题 24）。由此得出结论： \(u\) 必然是恒等映射（注意若 \(a\) 在 \(\mathbf{Q}\) 上是代数的，且 \(b\) 在 \(\mathbf{Q}\) 上是超越的，则 \(ab\) 在 \(\mathbf{Q}\) 上是超越的）.

\begin{enumerate}
\def\labelenumi{\arabic{enumi})}
\setcounter{enumi}{32}
\tightlist
\item
  设 \(\mathbf{K}\) 是一个交换域（特征任意， \(q\) )，并令 \(l\) 是一个素数。设 \(\mathbf{K}'\) 是 \(\mathbf{K}\) 的一个代数扩张，满足以下两个性质：
\end{enumerate}

\begin{enumerate}
\def\labelenumi{(\roman{enumi})}
\item
  \(\mathbf{K}[\mathbf{X}]\) 中每个次数不被 \(l\) 整除的非常数多项式在 \(\mathbf{K}'\) 中有一个根 ;
\item
  \(\mathbf{K}'[\mathbf{X}]\) 中每个次数等于 \(l\) 的多项式在 \(\mathbf{K}'\) 中有一个根 .
\end{enumerate}

\begin{enumerate}
\def\labelenumi{\alph{enumi})}
\item
  证明完全闭包 \(\mathbf{K}_1\) （ \(\mathbf{K}\) 的）包含于 \(\mathbf{K}'\) ，并且满足与 (i) 类似的条件。
\item
  K 的每一个次数不被 l 整除的代数扩张都同构于 \(K'\) 的某个子扩张（使用本原元素定理）。
\item
  设 L 是 \(K'\) 的一个包含 K 的子域，并设 M 为 L 的伽罗瓦扩张，其次数为 l 的幂。证明 M 同构于 \(K'\) 的一个子扩张（利用 l-群是幂零群这一事实）。
\end{enumerate}

\(d\) ）由此推出 \(\mathbf{K}'\) 是 \(\mathbf{K}\) 的一个代数闭包（若 \(\mathbf{F}\) 是 \(\mathbf{K}\) 的一个有限次数的（可分）扩张，考虑一个有限次数的伽罗瓦扩张 \(\mathrm{F}_1\supset \mathrm{F}\) ，以及 \(\mathrm{F}_1\) 在西罗 \(l\) -子群（ \(\mathrm{F}_1\) 在 \(\mathbf{K}\) 上的伽罗瓦群的）作用下的固定子域）.

\begin{enumerate}
\def\labelenumi{\arabic{enumi})}
\setcounter{enumi}{33}
\tightlist
\item
  设 E 是一个全序集。一个（戴德金）分割是指 E 的一个划分 (S, D)，其中 S 和 D 是两个非空子集，且 D 具有形式 M(A)，即 E 的某个非空子集 A 的上界集；D 有最小元 x 的充分必要条件是 \(D = \{x, \rightarrow\{, \text{所以 } S = \} \leftarrow, x\{.\)
\end{enumerate}

在全序阿贝尔群 G 中，一个分割 (S, D) 称为真分割，如果对所有 e \textgreater{} 0（在 G 中）存在 \(x' \in S\) 与 \(x'' \in D\) 使得 \(x'' - x' < e\) 。证明 (S, D) 是真分割当且仅当 D 在由 G 的上界子集组成的幺半群 \(\mathfrak{M}(G)\) （VI，第 35 页，习题 30）中可逆；所有分割都是真分割当且仅当 G 是阿基米德的（VI，第 35 页，习题 31）；因此 * 同构于 R 的一个子群（VI，第 36 页，习题 33）。 *

\begin{enumerate}
\def\labelenumi{\arabic{enumi})}
\setcounter{enumi}{34}
\tightlist
\item
  序域 E 的子域 K 在 E 中稠密，如果对于每个非空开区间 \(a, b \in E\) 存在 \(x \in K\) 使得 a \textless{} x \textless{} b. 证明此时 E 与 K 是可比较的（习题 11，a））。若取 E 为 R(X)，其中 R 是极大序域，且 X 相对于 R 是无穷大的（VI，第 43 页，习题 24），证明 E 与子域 \(K = R(X^{2})\) 是可比较的，但 K 在 E 中并不稠密。
\end{enumerate}

要使 K 在 E 中稠密，其充要条件是：对所有 \(x \in E\) ，由 \(K \cap J \leftarrow, x\) 与 \(K \cap [x, \rightarrow]\) 组成的 K 的分割是一个真分割（习题 34）。

\begin{enumerate}
\def\labelenumi{\arabic{enumi})}
\setcounter{enumi}{35}
\item
  设 K 是一个序域，并设 \(\tilde{K}\) 为所有使得分割 (S, D) 为真分割的上界集 D 所构成的集合；证明 \(\tilde{K}\) 是一个序域，其中加法取为幺半群 \(\mathfrak{M}(K)\) 的加法，而两个元素 \(D_{1}\) 与 \(D_{2}\) （均属于 \(\tilde{K}\) 且包含于 \(K_{+}\) ）的乘积定义为集合 \(\langle D_{1}D_{2}\rangle\) （VI，第 35 页，习题 30），其中 \(D_{1}D_{2}\) 表示由 \(x_{1}x_{2}\) （ \(x_{1}\in D_{1}\) ， \(x_{2}\in D_{2}\) ）组成的集合。证明：对于从 K 到序域 E 的稠密子域 F 上的任意递增同态 f，存在从 E 到 \(\tilde{K}\) 中的递增同态 g，使得 \(g\circ f\) 是 K 到 \(\tilde{K}\) 中的自然单射。
\item
  设 K 是一个序域，且设 n \textgreater{} 0 为整数，使得 K 的每个元素 x \textgreater{} 0 都等于一个幂 \(y^{n}\) ，其中 y \textgreater{} 0 。证明域 \(\tilde{K}\) 具有相同的性质。
\item
  假设域 K 是极大序的。证明该域 \(\tilde{\mathbf{K}}\) 也是极大序的。（可以通过反证法来论证：考虑 \(\tilde{\mathbf{K}}\) 的一个极大序的序代数扩张 E，并假设 \(\mathrm{E} \neq \tilde{\mathbf{K}}\) 。那么，在 \(\mathrm{E} \cap [\tilde{\mathbf{K}}\) 的元素中存在一个元素 \(w\) ，其极小多项式 \(f \in \tilde{\mathbf{K}}[\mathbf{X}]\) 的次数最小，为 \(n > 1\) . 证明存在一个区间 \(\mathbf{a}, b\) （在 E 中），使得 \(a, b \in \mathbf{K}\) 且 \(a < w < b\) ，以及 E 的一个元素 \(e > 0\) ，使得 \(f'(x) > e\) 对所有 \(x \in [\mathbf{a}, b]\) 成立，或 \(f'(x) < -e\) 对所有 \(x \in [\mathbf{a}, b]\) 成立。证明存在 K 中的 \(r > 0\) ，使得 \(f(a) \leqslant -r\) 且 \(r \leqslant f(b)\) ，或 \(f(a) \geqslant r\) 且 \(f(b) \leqslant -r\) 。用 K 的元素逼近 \(f\) 的系数，并证明这样得到一个多项式 \(g \in \mathbf{K}[\mathbf{X}]\) ，使得函数 \(x \mapsto g(x)\) 在 \([\mathbf{a}, \mathbf{b}]\) 上单调，并且在此区间内 K 的某一点处取零值，该点任意接近 \(w\) .)
\item
  设 K 是一个序域，E 是一个极大序的序代数扩张。证明域 E 的 K-自同构只有恒等映射。（首先证明从 E 到自身的每个 K-同态都是满射，并且是序域的同构；如果 \(x \in E\) 具有极小多项式 \(f \in K[X]\) ，那么 E 的自同构必然置换属于 E 的 f 的根。由此推出它固定 x。）
\item
  设 K 为序域。把形式幂级数域 \(E = K((X))\) 构造成一个序域，取元素 \textgreater0 为这样的形式幂级数 \(\sum_{n=-h}^{\infty} r_n X^n\) ：其
\end{enumerate}

最低次数的非零系数为 \textgreater0.

\begin{enumerate}
\def\labelenumi{\alph{enumi})}
\item
  证明如果 \(L = K(X)\) 是 E 的子域，由单未定元的有理函数组成，则 \(E = \tilde{L}\) （习题 36）在序域的同构意义下（注意对 L 中每个元素 a \textgreater{} 0，存在整数 n \textgreater{} 0 使得 \(0 < X^{n} < a\) ).
\item
  假设 K 是极大序的。证明诸域 \(\mathbf{K}((\mathbf{X}^{1/n}))\) 的并集 F（按与 E 相同的方式排序）是 E 的一个序代数扩张，并且是极大序的（利用 V 第 150 页习题 2）。
\item
  证明域 \(\widetilde{\mathbf{F}}\) （习题 36）是由形如 \(\sum_{r} \alpha_r X^r\) 的形式级数组成的域，其中 \(r\) 遍历
\end{enumerate}

有理数集 Q，其中 \(\alpha_{r} \in K\) ，且使得 \(\alpha_{r} \neq 0\) 的 r 所成的集合是一个递增序列，该序列要么有限，要么趋于 \(+\infty\) ；按照与 E 相同的方式对该域进行排序。特别地 \(\widetilde{F} \neq F\) .

\begin{enumerate}
\def\labelenumi{\arabic{enumi})}
\setcounter{enumi}{40}
\tightlist
\item
  设 K 为一个序域，E 为 K 上维数为 2 的向量空间。
\end{enumerate}

\begin{enumerate}
\def\labelenumi{\alph{enumi})}
\item
  证明两两不同的直线三元组（ \(D_{1}\) , \(D_{2}\) , \(D_{3}\) ，均经过 0）的集合在群 \(GL^{+}(E)\) 的作用下有两条轨道（考虑固定两条直线的子群）。
\item
  两两不同的半直线三元组（ \(\Delta_{1}\) , \(\Delta_{2}\) , \(\Delta_{3}\) ，以 0 为原点）的集合在 \(GL(E)\) 的作用下有 7 条轨道，而在 \(GL^{+}(E)\) 的作用下有 14 条轨道（同样的方法）。
\end{enumerate}

